\subsection{Multiplying Polynomials: FOIL Method} 
When we multiply two \makeblank[1.5in], there are \makeblanksmall possible products.
\begin{Example}Multiply $(A+B)(C+D)$ as a rectangle.
\begin{lpic}{}{4}
\end{lpic}
\end{Example}
\noindent The \makeblanksmall possible products are:\\
\begin{tabular}{cl}
\textbf{F}&Multiply the \textbf{F}irst terms of each binomial\\
\textbf{O}&Multiply the \textbf{O}uter pair of the product\\
\textbf{I}&Multiply the \textbf{I}nner pair of the product\\
\textbf{L}&Multiply the \textbf{L}ast terms of each binomial
\end{tabular}\\
Using this as a shortcut to multiply binomials is called the \term{FOIL method}. Remember, we can only use this to multiply two \makeblank[1.5in].
\begin{Example}Multiply $(2x-3)(x+5)$ as a rectangle and using the FOIL method.
\begin{lpic}{}{4}
\end{lpic}
\end{Example}
\noindent When we multiply two \makeblank[1in] binomials, we know we will combine the middle terms. With practice, we may skip writing this step.
\begin{Example}Multiply $(2x-3)(3x-2)$ using the FOIL method.
\begin{lpic}{}{1}
\end{lpic}
\end{Example}